\clearpage
\SourceLeaf{333}\pdfbookmark[2]{2. Zur Berechnung der Chorde aus der Zeit}{parabolisch-002}
\begin{center}
[2.]\par
\medskip
Zur Berechnung der Chorde \(\rho\) aus der Zeit \(=T\;[=t'-t]\), u[nd der] Summe der beiden Rad[ien] \(=r+r'\).
\end{center}
\[
6kT=(r+r'+\rho)^{\frac{3}{2}}-(r+r'-\rho)^{\frac{3}{2}}.
\]

Es sei
\[
\frac{\BareRoot{\left(1+\frac{\rho}{r+r'}\right)}-\BareRoot{\left(1-\frac{\rho}{r+r'}\right)}}{\BareRoot{8}}=\sin\psi
\]
\[
\frac{1-\BareRoot{\left(1-\frac{\rho\rho}{(r+r')^2}\right)}}{4}=\sin\psi^2
\]
\[
\frac{1+\BareRoot{\left(1-\frac{\rho\rho}{(r+r')^2}\right)}}{2}=\cos2\psi
\]
\[
\frac{\BareRoot{\left(1+\frac{\rho}{r+r'}\right)}+\BareRoot{\left(1-\frac{\rho}{r+r'}\right)}}{2}=\BareRoot{\cos2\psi}
\]
\[
\mathrm{[I]}\qquad \frac{\rho}{\BareRoot{8}\,.(r+r')}=\sin\psi\,.\BareRoot{\cos2\psi}.
\]

[Ferner ist, wenn vorübergehend \(t=\BareRoot{\left(1+\frac{\rho}{r+r'}\right)}\) und \(u=\BareRoot{\left(1-\frac{\rho}{r+r'}\right)}\) gesetzt wird,]
\[
\begin{aligned}
6kT&=(r+r')^{\frac{3}{2}}(t-u)(tt+tu+uu)\\
&=(r+r')^{\frac{3}{2}}(t-u)\left(3-\frac{(t-u)^2}{2}\right)
\end{aligned}
\]
[also, wegen] \(\frac{t-u}{\BareRoot{8}}=\sin\psi\):
\[
\mathrm{[II]}\qquad \frac{3}{\BareRoot{2}}\frac{kT}{(r+r')^{\frac{3}{2}}}=\sin3\psi.
\]
[Nachdem \(\psi\) aus Gleichung II gefunden ist, ergibt sich \(\rho\) aus Gleichung I.]
\begin{center}\rule{.3\linewidth}{.4pt}\end{center}
